% SPDX-License-Identifier: LicenseRef-SCROLL-NC-Research-1.0
% Generated from manuscript.md by pandoc 3.6.3 (-s --citeproc --bibliography research/paper/references.bib -t latex). Do not edit by hand. Not compiled locally: no TeX engine (DEC-002).
% Options for packages loaded elsewhere
\PassOptionsToPackage{unicode}{hyperref}
\PassOptionsToPackage{hyphens}{url}
\documentclass[
]{article}
\usepackage{xcolor}
\usepackage{amsmath,amssymb}
\setcounter{secnumdepth}{-\maxdimen} % remove section numbering
\usepackage{iftex}
\ifPDFTeX
  \usepackage[T1]{fontenc}
  \usepackage[utf8]{inputenc}
  \usepackage{textcomp} % provide euro and other symbols
\else % if luatex or xetex
  \usepackage{unicode-math} % this also loads fontspec
  \defaultfontfeatures{Scale=MatchLowercase}
  \defaultfontfeatures[\rmfamily]{Ligatures=TeX,Scale=1}
\fi
\usepackage{lmodern}
\ifPDFTeX\else
  % xetex/luatex font selection
\fi
% Use upquote if available, for straight quotes in verbatim environments
\IfFileExists{upquote.sty}{\usepackage{upquote}}{}
\IfFileExists{microtype.sty}{% use microtype if available
  \usepackage[]{microtype}
  \UseMicrotypeSet[protrusion]{basicmath} % disable protrusion for tt fonts
}{}
\makeatletter
\@ifundefined{KOMAClassName}{% if non-KOMA class
  \IfFileExists{parskip.sty}{%
    \usepackage{parskip}
  }{% else
    \setlength{\parindent}{0pt}
    \setlength{\parskip}{6pt plus 2pt minus 1pt}}
}{% if KOMA class
  \KOMAoptions{parskip=half}}
\makeatother
\usepackage{longtable,booktabs,array}
\usepackage{calc} % for calculating minipage widths
% Correct order of tables after \paragraph or \subparagraph
\usepackage{etoolbox}
\makeatletter
\patchcmd\longtable{\par}{\if@noskipsec\mbox{}\fi\par}{}{}
\makeatother
% Allow footnotes in longtable head/foot
\IfFileExists{footnotehyper.sty}{\usepackage{footnotehyper}}{\usepackage{footnote}}
\makesavenoteenv{longtable}
% definitions for citeproc citations
\NewDocumentCommand\citeproctext{}{}
\NewDocumentCommand\citeproc{mm}{%
  \begingroup\def\citeproctext{#2}\cite{#1}\endgroup}
\makeatletter
 % allow citations to break across lines
 \let\@cite@ofmt\@firstofone
 % avoid brackets around text for \cite:
 \def\@biblabel#1{}
 \def\@cite#1#2{{#1\if@tempswa , #2\fi}}
\makeatother
\newlength{\cslhangindent}
\setlength{\cslhangindent}{1.5em}
\newlength{\csllabelwidth}
\setlength{\csllabelwidth}{3em}
\newenvironment{CSLReferences}[2] % #1 hanging-indent, #2 entry-spacing
 {\begin{list}{}{%
  \setlength{\itemindent}{0pt}
  \setlength{\leftmargin}{0pt}
  \setlength{\parsep}{0pt}
  % turn on hanging indent if param 1 is 1
  \ifodd #1
   \setlength{\leftmargin}{\cslhangindent}
   \setlength{\itemindent}{-1\cslhangindent}
  \fi
  % set entry spacing
  \setlength{\itemsep}{#2\baselineskip}}}
 {\end{list}}
\usepackage{calc}
\newcommand{\CSLBlock}[1]{\hfill\break\parbox[t]{\linewidth}{\strut\ignorespaces#1\strut}}
\newcommand{\CSLLeftMargin}[1]{\parbox[t]{\csllabelwidth}{\strut#1\strut}}
\newcommand{\CSLRightInline}[1]{\parbox[t]{\linewidth - \csllabelwidth}{\strut#1\strut}}
\newcommand{\CSLIndent}[1]{\hspace{\cslhangindent}#1}
\setlength{\emergencystretch}{3em} % prevent overfull lines
\providecommand{\tightlist}{%
  \setlength{\itemsep}{0pt}\setlength{\parskip}{0pt}}
\usepackage{bookmark}
\IfFileExists{xurl.sty}{\usepackage{xurl}}{} % add URL line breaks if available
\urlstyle{same}
\hypersetup{
  pdftitle={Skeletons, prefetching and optimistic updates under controlled delays: time to usable content in a single-browser benchmark},
  pdfauthor={{[}OWNER PLACEHOLDER: author line. Value recorded at the freeze (C3): Mohamed Bechir Kefi, no affiliation. To be confirmed by the owner before release.{]}},
  hidelinks,
  pdfcreator={LaTeX via pandoc}}

\title{Skeletons, prefetching and optimistic updates under controlled
delays: time to usable content in a single-browser benchmark}
\author{{[}OWNER PLACEHOLDER: author line. Value recorded at the freeze
(C3): Mohamed Bechir Kefi, no affiliation. To be confirmed by the owner
before release.{]}}
\date{}

\begin{document}
\maketitle

\textbf{Status.} Independent research manuscript / preprint draft, not
peer-reviewed. Every number in this text comes from a committed result
file of this repository; the mapping is in
\texttt{research/publication/artifact\_manifest.json} and the audit in
\texttt{research/publication/evidence\_audit.md}. Two owner decisions
are still pending and are disclosed where they matter: the test reading
of DEC-011 (Section 4.5) and the owner-only facts marked as
placeholders.

\section{Abstract}\label{abstract}

Skeleton placeholders, intent-based prefetching and optimistic updates
are recommended to make interfaces feel faster, but feeling faster and
delivering usable content sooner are different quantities. We asked how
these strategies change the time from a click to the first frame showing
the final, correct and usable state (time to usable content, TTUC),
compared with a blank baseline. One mock interface was driven by
scripted input in Google Chrome 158.0.8127.38 on one laptop, with a
seeded latency and failure schedule shared by all arms. The protocol was
frozen before the confirmatory runs, which covered 6,000 interactions in
the main experiments and 3,600 in a sensitivity experiment; the seed was
the independent unit. Prefetching with a 150 ms lead reduced TTUC at
every delay, by 50.0 to 160.5 ms (Holm-adjusted p at most 0.00053),
close to the mechanical expectation. An optimistic update reached the
final state about one delay earlier on every successful save (for
example -998.6 ms, 95\% CI -1012.2 to -972.6, at 1000 ms) and showed no
detectable TTUC difference on failed saves, while displaying a state the
server never held for about the whole delay. A skeleton was equivalent
to the blank baseline within 17 ms at delays of 100 ms and more; at 50
ms it was 16.7 ms faster under the main paint-time method, but this
equivalence verdict changed under both alternative methods and is not
robust. The results describe timing in one browser and one synthetic
page; they make no claim about perceived speed.

\section{Keywords}\label{keywords}

frontend performance; loading strategies; skeleton screens; prefetching;
optimistic updates; reproducible benchmarking

\section{1. Introduction}\label{introduction}

Interfaces that wait for a server must decide what to show during the
wait. Common strategies show nothing new until the response arrives,
show a placeholder with the final layout (a skeleton), start the request
earlier on an intent signal such as a hover (prefetching), or show the
expected result immediately and correct it if the server disagrees (an
optimistic update). Practitioner guidance recommends skeletons and
reserved space to reduce layout shift (Osmani and Pollard, n.d.;
Tankala, n.d.), and browser standards define how the latency from an
input to the next paint is measured (Mocny, n.d.a; Wagner and Pollard,
n.d.). These strategies act on different quantities. A skeleton changes
what is shown first, but it cannot make the data arrive sooner; a
prefetch can; an optimistic update shows the final state at once on
success and a false state on failure.

The project's working question, set in its master prompt, is: how do
loading strategies affect measurable interface responsiveness under
controlled delays? We operationalise it as the time from a click to the
first frame showing the requested content in its final, correct and
usable state (TTUC), and we keep the time to first feedback (TTFF) as a
secondary metric. The claims boundary is strict: objective timing is not
perceived speed, and nothing here is a statement about human perception.

The gap addressed is narrow. The short, non-systematic search of the
project (Section 2) found no controlled objective comparison of these
strategies against a blank baseline under identical seeded delays; this
is not evidence that none exists.

Hypotheses (frozen protocol v1.0, Section 4): for each non-baseline arm
and condition, H0 is a zero trial-paired TTUC difference against the
blank baseline, and H1 is two-sided. The pre-stated directional
expectations were: no TTUC benefit for SKELETON (equivalence within 17
ms), a benefit of about -min(h, L) for PREFETCH, and a large benefit for
OPTIMISTIC at p = 0 that shrinks as the failure probability rises, with
growing wrong-state exposure.

Contributions:

\begin{enumerate}
\def\labelenumi{\arabic{enumi}.}
\tightlist
\item
  A frozen, reproducible benchmark in which one mock page, one seeded
  schedule and one scripted input sequence are shared by four loading
  strategies, with every raw trial, manifest and checksum committed.
\item
  Confirmatory estimates for 23 primary contrasts (8 on the read path,
  15 on the write path) with seed-level uncertainty and multiplicity
  control, including null results.
\item
  A pre-registered sensitivity experiment over paint-time method, CPU
  throttling, prefetch lead and headed rendering, which shows which
  verdicts depend on how a frame is timed.
\item
  An exploratory post hoc decomposition of the one frame-scale effect,
  from committed traces only, that locates it in frame alignment rather
  than in rendering cost.
\end{enumerate}

\section{2. Related work}\label{related-work}

This section is a scoping summary, not a systematic review. Every source
was opened on 2026-11-24 through a tool that returns a model-written
summary of the page rather than its full text
(\texttt{research/literature\_review.md}). Claims below are limited to
what those summaries stated; a full-text reading of every source a claim
depends on is scheduled for M6 (owner, 2026-11-25, ``Approve all''
(post-M4), item 7).

\textbf{Measurement standards.} The Event Timing API defines an event's
duration from its timestamp to the next rendering update after dispatch,
rounded to 8 ms, with an observer threshold of at least 16 ms (Mocny,
n.d.a). Interaction to Next Paint aggregates such durations into a
page-level responsiveness metric (Wagner and Pollard, n.d.). The Layout
Instability API defines layout-shift values and excludes shifts shortly
after discrete input (Mocny, n.d.b; Mihajlija and Walton, n.d.). Largest
Contentful Paint stops reporting at the first input, so it cannot time
content that appears after a click (MDN Web Docs, n.d.). These sources
motivate the custom TTUC marker used here and explain why Event Timing
could only be descriptive: most of our interactions fell below its 16 ms
threshold.

\textbf{Strategy guidance.} Practitioner guidance recommends reserving
space with a placeholder or skeleton to avoid late layout shifts (Osmani
and Pollard, n.d.) and discusses when skeleton screens are appropriate
(Tankala, n.d.). Response-time limits of 0.1, 1 and 10 s (Nielsen 1993)
were used only to choose latency levels. The React documentation of
optimistic state describes the rollback semantics our optimistic arm
imitates (React, n.d.).

\textbf{Perception studies.} A study of skeleton screens and perceived
speed exists (Mejtoft, Långström, and Söderström 2018); only its
bibliographic metadata was verified and its findings are not used.
Perception is outside the scope of this work.

\textbf{Relation to this work.} The study does not replicate any prior
experiment. It extends standard browser instrumentation with a
usable-content marker and applies it to a controlled comparison; its
contribution is methodological and descriptive about one page in one
browser.

\section{3. Background and formal
model}\label{background-and-formal-model}

Let \texttt{t\_in} be the click event's timestamp on the page clock. For
a trial, the mock transport returns after a scheduled delay
\texttt{d\ =\ L\ (1\ +\ j)}, with \texttt{j} uniform in {[}-0.1, 0.1{]}
and drawn from a seeded stream, and the write path fails with
probability \texttt{p}. The same seed yields the same delays and
outcomes in every arm (common random numbers).

TTUC is \texttt{t\_final\ -\ t\_in}, where \texttt{t\_final} is the time
of the first frame from which the target region shows the final,
server-consistent, fully rendered and actionable state and keeps it
until the trial ends. On a failed save, the final state is the unchanged
old state together with the error message. The main paint-time method
takes \texttt{performance.now()} inside the second
\texttt{requestAnimationFrame} callback after the DOM commit. Two
alternatives are recorded by the frozen code: the timestamp argument of
the first \texttt{requestAnimationFrame} callback (rAF1, the start of
the frame that will carry the change) and, on the read path, the Element
Timing \texttt{renderTime} of the content.

A reference model gives each arm's ideal timing with zero overhead: for
BLANK and SKELETON on the read path, \texttt{TTUC\ \textgreater{}=\ d};
for PREFETCH, \texttt{TTUC\ \textgreater{}=\ max(0,\ d\ -\ h)}, where
\texttt{h} is the hover lead; for OPTIMISTIC,
\texttt{TTUC\ \textgreater{}=\ 0} on success and
\texttt{\textgreater{}=\ d} on failure. The excess over this bound is
rendering and scheduling overhead.

For arm \texttt{s} and cell \texttt{c}, the estimand (amendment A1) is

\texttt{Delta\_s(c)\ =\ median\ over\ seeds\ of\ {[}\ median\ over\ k\ of\ (\ TTUC\_s,k\ -\ TTUC\_BLANK,k\ )\ {]}},

where trial \texttt{k} in both arms shares its scheduled delay and
outcome. On the write path with \texttt{p\ \textgreater{}\ 0}, success
and failure trials form separate primary strata (amendment A2); the
mixed estimand over all trials is secondary.

\section{4. Methods}\label{methods}

\subsection{4.1 Design and protocol
registration}\label{design-and-protocol-registration}

All confirmatory work followed \texttt{research/protocol.md} v1.0,
frozen by the owner on 2026-11-24 (SHA-256
\texttt{a87ae45e049c16ed42abef0ab853a0e3fcf28daec9d5482d11bbb1dfb365f004}),
with per-experiment protocols and configurations whose hashes are
recorded in \texttt{research/freeze.json}. A run guard refused any
confirmatory run whose protocol or configuration differed from the
frozen one. Six amendments (A1 to A6) were approved before the freeze,
after an exploratory pre-freeze run whose values of the contrasts were
seen (freeze record, C4 disclosure). EXP-001 (instrumentation
validation) is exploratory by design. EXP-002 (read path), EXP-003
(write path) and EXP-004 (sensitivity) are confirmatory. The analyses of
Section 5.5 are exploratory and post hoc.

\subsection{4.2 System under test and loading
strategies}\label{system-under-test-and-loading-strategies}

One locally served page shows a list of items and a detail panel,
implemented once with the strategy as the only switch
(\texttt{src/web/}). On the read path, a click on an item requests its
detail: BLANK leaves the region empty until the response, SKELETON
renders a placeholder with the final box size at input, and PREFETCH
starts the request on a scripted hover 150 ms before the click. On the
write path, a click on a save toggle sends a mutation: BLANK keeps the
old state until the server confirms, and OPTIMISTIC shows the new state
at once and rolls it back with an error message on failure. The input
script is identical in every arm, and only PREFETCH reacts to the hover.

The page ran in the installed Google Chrome 158.0.8127.38, driven by
Playwright 1.65.0 through \texttt{channel\ =\ "chrome"} (amendment A5),
headless at CPU 1x except in EXP-004 S2 and S4. The machine was an Apple
M1 Pro (10 cores, 16 GiB) under macOS 15.3.1 on mains power.

\subsection{4.3 Latency and failure
injection}\label{latency-and-failure-injection}

Read latencies were L in \{50, 100, 300, 1000\} ms and write latencies L
in \{100, 300, 1000\} ms, with p in \{0, 0.1, 0.3\}. Each session is a
fresh browser context with K = 10 scripted interactions of one path;
arms are ordered within each seed block by a seeded permutation.
Confirmatory seeds were 1 to 20 (EXP-002, EXP-003) and 1 to 10
(EXP-004). With K = 10 the schedule fixes, before data, that the failure
stratum has 12 seeds at p = 0.1 and 18 seeds at p = 0.3 (9 seeds in
EXP-004).

\subsection{4.4 Metrics and
instrumentation}\label{metrics-and-instrumentation}

The primary metric is TTUC (Section 3), in milliseconds, lower is
better. Secondary metrics are TTFF, layout shift (all entries and
CLS-eligible entries), wrong-state exposure on the write path (split
into false state, a state the server never held, and stale state, a
superseded state), rollback correctness, and Event Timing duration,
which is descriptive only (amendment A4). Exclusions were allowed only
for harness failure and for the hover-lead control (amendment A3: a read
trial is excluded when the measured hover-to-click lead deviates from h
by more than 25 ms). A cell with more than 5\% of its trials excluded is
a failed cell. The stop rule required rollback correctness of 1.0 and no
BLANK trial below its injected delay.

\subsection{4.5 Statistical analysis}\label{statistical-analysis}

The seed is the independent unit. Each contrast is summarised by the
median of its per-seed paired differences, with a 95\% percentile
bootstrap interval over seeds (10,000 resamples, fixed seed). The
hypothesis test is the two-sided Wilcoxon signed-rank test on the
per-seed values, with Holm correction over each experiment's primary
family (8 contrasts in EXP-002, 15 in EXP-003). SKELETON equivalence
uses two one-sided tests through the 90\% bootstrap interval with a
margin of 17 ms. Non-significance is never reported as equivalence.

The frozen protocol names the bootstrap interval, the Wilcoxon test and
the Holm correction without stating which one decides H0. Before any
confirmatory output existed, DEC-011 fixed the reading ``H0 is rejected
when the Holm-adjusted Wilcoxon p is below 0.05'' and required the
alternative reading, ``the unadjusted 95\% interval excludes 0'', to be
reported beside it. \textbf{DEC-011 awaits the owner's review and is not
decided here.} Both readings are shown in Table 1 and Table 2; they
differ for exactly one of the 23 primary contrasts (Section 5.2).

EXP-004 applies the frozen robustness rule: a verdict is robust only if
the sign of the point estimate is the same at every available factor
level; for SKELETON the rule uses the equivalence verdict. DEC-011 item
5 also fixed which alternative paint methods form level S1 (both rAF1
and Element Timing); this too is part of the pending owner review, and
Section 5.4 shows the effect of the narrower reading.

\subsection{4.6 Reproducibility}\label{reproducibility}

Commands are in each experiment report's reproduction appendix. Raw
results are committed under \texttt{results/raw/} with SHA-256
manifests; all 110 listed files matched their checksums in the M4 audit.
Tables and figures are produced by the frozen generators
\texttt{python\ -m\ scroll\ report} and
\texttt{python\ -m\ scroll.inference}, and the exploratory analyses by
\texttt{scripts/m4\_reanalysis.py}. Confirmatory runs used code revision
\texttt{b5a72c3}, whose run path is identical to the freeze revision
\texttt{b4bd25b}.

\section{5. Results}\label{results}

\subsection{5.1 EXP-001: instrumentation validation
(exploratory)}\label{exp-001-instrumentation-validation-exploratory}

EXP-001 was run twice on exploratory seeds 1001 to 1005, which are
disjoint from the confirmatory seeds: once before the freeze on a loaded
machine (\texttt{EXP-001-20261124T161058Z}, load average 25.31 at the
start), and once as the metric validation in the quiet window before the
confirmatory runs (\texttt{EXP-001-20261125T065102Z}). Both runs had 70
of 70 sessions and 350 interactions. All four instrumentation checks
passed in both: no BLANK trial below its injected delay (minimum excess
2.5 and 3.2 ms), cell median excess over the reference model of 11.5 to
23.8 ms and 12.0 to 27.0 ms, against a planning threshold of 50 ms,
rollback correct in every write trial, and prefetched responses used in
every PREFETCH trial. The second run had one hover-lead exclusion. The
pre-freeze run motivated amendments A1 to A3 and A6; its values are not
evidence. The protocol said EXP-001 would choose the alternative paint
method for S1; no choice was recorded (DEC-012).

\subsection{5.2 EXP-002: read path}\label{exp-002-read-path}

EXP-002 (\texttt{EXP-002-20261125T070315Z}) completed 240 of 240
sessions in 3,887.4 s; 2,387 of 2,400 trials were analysed after 13
hover-lead exclusions (at most 2 of 200 in any cell and arm), and no
cell failed. No BLANK trial had TTUC below its injected delay (0 of
796), and the final state was correct in every trial.

Table 1. EXP-002 primary contrasts. Median trial-paired TTUC difference
against BLANK, ms; 20 seeds per contrast; 95\% percentile bootstrap
interval; two-sided Wilcoxon with Holm over 8. Reading A is DEC-011
(Holm-adjusted Wilcoxon); reading B is the unadjusted 95\% interval.
Source:
\texttt{results/raw/EXP-002/EXP-002-20261125T070315Z/report/inference.json}.

\begin{longtable}[]{@{}
  >{\raggedright\arraybackslash}p{(\linewidth - 14\tabcolsep) * \real{0.1250}}
  >{\raggedright\arraybackslash}p{(\linewidth - 14\tabcolsep) * \real{0.1250}}
  >{\raggedright\arraybackslash}p{(\linewidth - 14\tabcolsep) * \real{0.1250}}
  >{\raggedright\arraybackslash}p{(\linewidth - 14\tabcolsep) * \real{0.1250}}
  >{\raggedright\arraybackslash}p{(\linewidth - 14\tabcolsep) * \real{0.1250}}
  >{\raggedright\arraybackslash}p{(\linewidth - 14\tabcolsep) * \real{0.1250}}
  >{\raggedright\arraybackslash}p{(\linewidth - 14\tabcolsep) * \real{0.1250}}
  >{\raggedright\arraybackslash}p{(\linewidth - 14\tabcolsep) * \real{0.1250}}@{}}
\toprule\noalign{}
\begin{minipage}[b]{\linewidth}\raggedright
L (ms)
\end{minipage} & \begin{minipage}[b]{\linewidth}\raggedright
Contrast
\end{minipage} & \begin{minipage}[b]{\linewidth}\raggedright
Median (ms)
\end{minipage} & \begin{minipage}[b]{\linewidth}\raggedright
95\% CI (ms)
\end{minipage} & \begin{minipage}[b]{\linewidth}\raggedright
Holm p
\end{minipage} & \begin{minipage}[b]{\linewidth}\raggedright
Reading A
\end{minipage} & \begin{minipage}[b]{\linewidth}\raggedright
Reading B
\end{minipage} & \begin{minipage}[b]{\linewidth}\raggedright
TOST 90\% CI (ms), equivalence
\end{minipage} \\
\midrule\noalign{}
\endhead
\bottomrule\noalign{}
\endlastfoot
50 & PREFETCH & -50.0 & {[}-51.2, -48.9{]} & 0.00053 & reject & excludes
0 & \\
50 & SKELETON & -16.7 & {[}-17.1, -16.0{]} & 0.00053 & reject & excludes
0 & {[}-17.05, -16.1{]}, not shown \\
100 & PREFETCH & -84.9 & {[}-86.3, -84.1{]} & 1.5e-05 & reject &
excludes 0 & \\
100 & SKELETON & -0.6 & {[}-1.9, 0.1{]} & 0.27 & not rejected & includes
0 & {[}-1.6, 0.0{]}, equivalent \\
300 & PREFETCH & -160.5 & {[}-163.5, -157.1{]} & 1.5e-05 & reject &
excludes 0 & \\
300 & SKELETON & -1.0 & {[}-1.1, -0.4{]} & 0.27 & not rejected &
\textbf{excludes 0} & {[}-1.1, -0.6{]}, equivalent \\
1000 & PREFETCH & -157.8 & {[}-159.1, -155.4{]} & 0.00053 & reject &
excludes 0 & \\
1000 & SKELETON & 0.0 & {[}-0.9, 0.9{]} & 1.0 & not rejected & includes
0 & {[}-0.7, 0.8{]}, equivalent \\
\end{longtable}

\textbf{PREFETCH.} Prefetching reduced TTUC at every delay. The
estimates were within 17 ms of the mechanical expectation -min(h, L):
deviations 0.0, 15.1, -10.5 and -7.8 ms at L = 50, 100, 300 and 1000 ms,
so the frozen contradiction rule was not met. At L of 50 and 100 ms the
residual wait is zero by construction, and TTUC equals the rendering
cost (cell median 27.7 ms); the deviation at L = 100 ms arises because
BLANK's own overhead there (median excess 16.3 ms) is smaller than
PREFETCH's.

\textbf{SKELETON at L of 100 ms and more.} The skeleton was equivalent
to BLANK within 17 ms, with estimates of -1.0 to 0.0 ms, while it
lowered TTFF from about the full delay to about 28 ms (cell medians 27.5
to 28.0 ms). The two test readings disagree at L = 300 ms: the
unadjusted interval {[}-1.1, -0.4{]} ms excludes 0, the Holm-adjusted
Wilcoxon test does not reject. The difference is about 1 ms, and neither
reading changes the equivalence verdict.

\textbf{SKELETON at L = 50 ms.} The skeleton reached usable content 16.7
ms earlier than BLANK under the main paint method, H0 was rejected, and
equivalence was not shown because the lower 90\% bound, -17.05 ms, lies
just beyond the margin. Section 5.4 shows that this verdict is not
robust.

Layout shift was identical across arms in most cells (median 0.035 in
all entries); the CLS-eligible median was 0.035 for BLANK and PREFETCH
at L = 1000 ms and 0 elsewhere, a design consequence of BLANK's
collapsing empty region (DEC-007). Event Timing was below its 16 ms
threshold in 71 to 79\% of BLANK trials. Figures:
\texttt{results/raw/EXP-002/EXP-002-20261125T070315Z/report/figures/fig-ttuc.svg}
(TTUC per trial by cell and arm, log scale) and \texttt{fig-paired.svg}
(per-seed paired differences with 95\% intervals).

\subsection{5.3 EXP-003: write path}\label{exp-003-write-path}

EXP-003 (\texttt{EXP-003-20261125T080838Z}) completed 360 of 360
sessions in 6,318.7 s, with 3,600 of 3,600 trials analysed, no exclusion
and rollback correctness 1.0; the stop rule was not triggered.

Table 2. EXP-003 primary contrasts, OPTIMISTIC minus BLANK, ms. Seeds
per stratum fixed by the schedule; Holm over 15. Reading A and reading B
as in Table 1; they agree for all 15 contrasts. Source:
\texttt{results/raw/EXP-003/EXP-003-20261125T080838Z/report/inference.json}.

\begin{longtable}[]{@{}
  >{\raggedright\arraybackslash}p{(\linewidth - 16\tabcolsep) * \real{0.1111}}
  >{\raggedright\arraybackslash}p{(\linewidth - 16\tabcolsep) * \real{0.1111}}
  >{\raggedright\arraybackslash}p{(\linewidth - 16\tabcolsep) * \real{0.1111}}
  >{\raggedright\arraybackslash}p{(\linewidth - 16\tabcolsep) * \real{0.1111}}
  >{\raggedright\arraybackslash}p{(\linewidth - 16\tabcolsep) * \real{0.1111}}
  >{\raggedright\arraybackslash}p{(\linewidth - 16\tabcolsep) * \real{0.1111}}
  >{\raggedright\arraybackslash}p{(\linewidth - 16\tabcolsep) * \real{0.1111}}
  >{\raggedright\arraybackslash}p{(\linewidth - 16\tabcolsep) * \real{0.1111}}
  >{\raggedright\arraybackslash}p{(\linewidth - 16\tabcolsep) * \real{0.1111}}@{}}
\toprule\noalign{}
\begin{minipage}[b]{\linewidth}\raggedright
L (ms)
\end{minipage} & \begin{minipage}[b]{\linewidth}\raggedright
p
\end{minipage} & \begin{minipage}[b]{\linewidth}\raggedright
Stratum
\end{minipage} & \begin{minipage}[b]{\linewidth}\raggedright
Seeds
\end{minipage} & \begin{minipage}[b]{\linewidth}\raggedright
Median (ms)
\end{minipage} & \begin{minipage}[b]{\linewidth}\raggedright
95\% CI (ms)
\end{minipage} & \begin{minipage}[b]{\linewidth}\raggedright
Holm p
\end{minipage} & \begin{minipage}[b]{\linewidth}\raggedright
Reading A
\end{minipage} & \begin{minipage}[b]{\linewidth}\raggedright
Reading B
\end{minipage} \\
\midrule\noalign{}
\endhead
\bottomrule\noalign{}
\endlastfoot
100 & 0 & success & 20 & -85.2 & {[}-88.0, -84.3{]} & 0.00079 & reject &
excludes 0 \\
100 & 0.1 & success & 20 & -85.3 & {[}-86.6, -84.8{]} & 2.9e-05 & reject
& excludes 0 \\
100 & 0.1 & failure & 12 & 0.5 & {[}-1.5, 1.7{]} & 1.0 & not rejected &
includes 0 \\
100 & 0.3 & success & 20 & -85.3 & {[}-88.0, -83.9{]} & 0.00079 & reject
& excludes 0 \\
100 & 0.3 & failure & 18 & 2.7 & {[}-1.0, 5.4{]} & 0.37 & not rejected &
includes 0 \\
300 & 0 & success & 20 & -288.5 & {[}-297.7, -282.3{]} & 0.00079 &
reject & excludes 0 \\
300 & 0.1 & success & 20 & -289.3 & {[}-297.6, -283.5{]} & 2.9e-05 &
reject & excludes 0 \\
300 & 0.1 & failure & 12 & -0.8 & {[}-3.0, 4.3{]} & 1.0 & not rejected &
includes 0 \\
300 & 0.3 & success & 20 & -290.4 & {[}-297.4, -284.1{]} & 2.9e-05 &
reject & excludes 0 \\
300 & 0.3 & failure & 18 & -0.1 & {[}-2.6, 1.5{]} & 1.0 & not rejected &
includes 0 \\
1000 & 0 & success & 20 & -998.6 & {[}-1012.2, -972.6{]} & 2.9e-05 &
reject & excludes 0 \\
1000 & 0.1 & success & 20 & -995.9 & {[}-1015.8, -967.3{]} & 2.9e-05 &
reject & excludes 0 \\
1000 & 0.1 & failure & 12 & -1.6 & {[}-7.1, 0.1{]} & 0.46 & not rejected
& includes 0 \\
1000 & 0.3 & success & 20 & -990.7 & {[}-1016.6, -961.7{]} & 2.9e-05 &
reject & excludes 0 \\
1000 & 0.3 & failure & 18 & 0.6 & {[}-1.4, 3.4{]} & 1.0 & not rejected &
includes 0 \\
\end{longtable}

On success, the optimistic update showed the final state about one delay
earlier; OPTIMISTIC cell medians of TTUC were 27.0 to 29.7 ms against
111.6 to 1024.8 ms for BLANK. On failure, no stratum rejected H0, with
estimates from -1.6 to 2.7 ms; no equivalence test was registered for
these strata, so they are not equivalence claims. The secondary mixed
estimand decreased strictly in magnitude from p = 0 to 0.1 to 0.3 at
each delay (for example -998.6, -974.1 and -951.5 ms at L = 1000 ms), as
expected, so the frozen trend check found no contradiction.

The cost of the optimistic update is visible in wrong-state exposure
rather than in TTUC. In failure trials (83 per delay and arm over both p
levels), OPTIMISTIC displayed a false state for a median of 83.4, 283.9
and 999.1 ms at L = 100, 300 and 1000 ms, and BLANK displayed none.
Layout shift was 0 in every write cell. Figures:
\texttt{results/raw/EXP-003/EXP-003-20261125T080838Z/report/figures/fig-ttuc.svg}
and \texttt{fig-paired.svg}.

\subsection{5.4 EXP-004: sensitivity}\label{exp-004-sensitivity}

EXP-004 ran S2 (CPU 4x, 140 sessions), S3 (h = 0 and 600 ms, 80
sessions, 4 hover-lead exclusions) and S4 (headed, 140 sessions) in
full, 3,596 of 3,600 trials analysed, with no failed cell; S1 recomputed
every EXP-002 and EXP-003 contrast with the alternative paint methods.
Under the frozen rule, 19 of 23 primary contrasts were robust, including
every PREFETCH and every OPTIMISTIC success contrast and the SKELETON
equivalence at L of 100 ms and more.

Table 3. The four contrasts that the frozen rule does not call robust.
Median trial-paired TTUC difference, ms, with 95\% interval. Source:
\texttt{results/raw/EXP-004/sensitivity/sensitivity.json}.

\begin{longtable}[]{@{}
  >{\raggedright\arraybackslash}p{(\linewidth - 12\tabcolsep) * \real{0.1429}}
  >{\raggedright\arraybackslash}p{(\linewidth - 12\tabcolsep) * \real{0.1429}}
  >{\raggedright\arraybackslash}p{(\linewidth - 12\tabcolsep) * \real{0.1429}}
  >{\raggedright\arraybackslash}p{(\linewidth - 12\tabcolsep) * \real{0.1429}}
  >{\raggedright\arraybackslash}p{(\linewidth - 12\tabcolsep) * \real{0.1429}}
  >{\raggedright\arraybackslash}p{(\linewidth - 12\tabcolsep) * \real{0.1429}}
  >{\raggedright\arraybackslash}p{(\linewidth - 12\tabcolsep) * \real{0.1429}}@{}}
\toprule\noalign{}
\begin{minipage}[b]{\linewidth}\raggedright
Contrast
\end{minipage} & \begin{minipage}[b]{\linewidth}\raggedright
Main
\end{minipage} & \begin{minipage}[b]{\linewidth}\raggedright
S1 rAF1
\end{minipage} & \begin{minipage}[b]{\linewidth}\raggedright
S1 Element
\end{minipage} & \begin{minipage}[b]{\linewidth}\raggedright
S2 CPU 4x
\end{minipage} & \begin{minipage}[b]{\linewidth}\raggedright
S4 headed
\end{minipage} & \begin{minipage}[b]{\linewidth}\raggedright
Frozen verdict
\end{minipage} \\
\midrule\noalign{}
\endhead
\bottomrule\noalign{}
\endlastfoot
read L=50 SKELETON & -16.7 {[}-17.1, -16.0{]}, eq no & -15.3 {[}-16.0,
-14.8{]}, eq yes & -9.1 {[}-9.8, -8.3{]}, eq yes & & &
method-dependent \\
write L=300 p=0.1 failure & -0.8 {[}-3.0, 4.3{]} & 0.7 {[}-1.9, 4.3{]} &
& & & method-dependent \\
write L=300 p=0.3 failure & -0.1 {[}-2.6, 1.5{]} & 0.4 {[}-1.4, 1.7{]} &
& -0.3 {[}-0.9, 2.0{]} & 0.4 {[}-2.7, 0.9{]} & method- and
condition-dependent \\
write L=1000 p=0.3 failure & 0.6 {[}-1.4, 3.4{]} & 1.2 {[}-0.8, 2.1{]} &
& -1.3 {[}-4.2, 0.9{]} & 0.0 {[}-1.3, 2.8{]} & condition-dependent \\
\end{longtable}

\textbf{The 50 ms skeleton effect is not robust.} Its equivalence
verdict changes under both alternative paint methods, so under the
frozen rule the result is method-dependent. The direction is the same
under all three methods, but the size is not (-16.7, -15.3 and -9.1 ms),
and whether it exceeds the equivalence margin depends on the method. We
therefore do not report a skeleton benefit at 50 ms.

\textbf{The three failure strata.} Their main estimates lie within 1 ms
of zero, none rejects H0, and the frozen sign rule flags them because
the sign of a near-zero estimate flips. Under the exploratory interval
reading of Section 5.5, every level's 95\% interval includes 0 for all
six failure strata, which is the pattern expected of a null contrast.
The frozen verdicts stand as reported.

\textbf{Reading of level S1 (pending owner review).} The frozen EXP-004
protocol names the S1 alternative as ``Element Timing or trace''.
DEC-011 counted rAF1 as well, which made S1 available on the write path.
Under the narrower reading (Element Timing only, which exists on the
read path only; X6 in Section 5.5), the read-path verdicts are
unchanged, but the nine write contrasts at L = 100 ms or p = 0.1 have no
sensitivity level at all, because S2 to S4 cover only L in \{300, 1000\}
ms and p in \{0, 0.3\}. The count becomes 11 robust, 1 method-dependent
(the 50 ms skeleton), 2 condition-dependent (write L = 300 and 1000 ms,
p = 0.3, failure) and 9 not assessed. The skeleton result at 50 ms is
method-dependent under both readings.

\textbf{Shifts within robust contrasts.} Under CPU 4x the prefetch
benefit shrank by 5.9 to 7.6 ms and the optimistic success benefit grew
by 5.9 to 13.7 ms; headed rendering enlarged the optimistic success
benefit by 15.3 to 27.0 ms (10 seeds each; differences of point
estimates, not tests). The headed shift has an untested competing
explanation: in S4 the median second-frame gap was 8.6 to 9.0 ms for
OPTIMISTIC but 16.45 to 16.75 ms for BLANK on the write path
(exploratory X2), which suggests that the display ran a faster frame
cadence shortly after input; the refresh rate was not recorded, and the
S4 run also took place one to three hours after the main runs. In S3 the
prefetch benefit scaled with the lead: -293.9 and -609.6 ms at h = 600
ms (L = 300 and 1000 ms), and -1.9 and -1.7 ms at h = 0. TTFF
sensitivity could not be computed, because the frozen code records
alternative paint times only for the usable-content frame. Figure:
\texttt{results/raw/EXP-004/sensitivity/figures/fig-sensitivity.svg}; it
shows only the 10 contrasts of the reduced matrix, so the 50 ms skeleton
contrast and the contrasts checked under S1 only appear in Table 3 and
the EXP-004 report, not in the figure.

\subsection{5.5 Exploratory re-analysis of committed traces (M4, post
hoc)}\label{exploratory-re-analysis-of-committed-traces-m4-post-hoc}

The following analyses were run on 2026-11-25 by
\texttt{scripts/m4\_reanalysis.py} (code revision and input checksums in
the output manifest), on the committed \texttt{results.json} files only,
with no browser run. They are exploratory, post hoc and never replace a
frozen verdict. Outputs: \texttt{results/exploratory/M4-REANALYSIS/}.

\textbf{X1, test readings.} Tabulating both readings over the 23 primary
contrasts confirms that they disagree only for EXP-002 SKELETON at L =
300 ms (\texttt{x1\_test\_readings.csv}).

\textbf{X2, where the 50 ms skeleton effect sits.} The double-rAF TTUC
can be split into the first-frame time relative to the reference delay
and the gap to the second frame. At L = 50 ms the second-frame gap was
the same in both arms (median 18.9 ms for BLANK, 18.8 ms for SKELETON),
whereas the first frame began 8.1 ms after the reference delay for BLANK
and 7.0 ms before it for SKELETON. In other words, for SKELETON the
frame that carries the content was already under way when the content
was committed (83\% of trials), and for BLANK it almost never was (1\%).
At L of 100 ms and more, BLANK behaved like SKELETON (first frame before
the reference in 72 to 79\% of read trials). At the trial level the
main-method difference was concentrated, not driven by a few outliers:
median -16.5 ms, interquartile range -19.3 to -14.2 ms, and 80.5\% of
the 200 trial pairs between -25 and -8 ms. Under Element Timing the
trial-level median was -9.3 ms. These observations locate the effect in
frame alignment, the competing explanation that the EXP-002 protocol
named (``frame alignment could create artificial 8 to 17 ms steps''),
not in a rendering cost. Why the renderer was idle for BLANK only at 50
ms is not tested (\texttt{x2\_frame\_decomposition.csv},
\texttt{x2\_skeleton\_l50\_trial\_distribution.csv}).

\textbf{X2, frame interval against the margin.} In the headless runs the
median second-frame gap, an estimate of the frame interval, was 18.0 to
18.9 ms in every cell, larger than the 17 ms margin that the protocol
justified as one 60 Hz frame. A difference of one frame under the main
method therefore cannot pass the equivalence test.

\textbf{X2, the rAF1 method.} The rAF1 time precedes the reference delay
in most trials (for example 72 to 79\% of BLANK read trials at L of 100
ms and more), because it is the start of the frame rather than a paint
time. It is a lower, optimistic estimate, as the code documents, and the
instrumentation check ``no BLANK trial below the injected delay'' holds
only for the main method (0\% of trials).

\textbf{X3, hover-lead exclusions.} Keeping the 13 trials excluded by
the A3 rule changes no EXP-002 verdict; the largest change of a point
estimate is 0.2 ms (\texttt{x3\_a3\_inclusion\_exp002.csv}).

\textbf{X4, interval reading of robustness.} For every non-SKELETON
contrast, we classified each factor level by whether its 95\% interval
excludes 0. Twelve of the 13 PREFETCH and OPTIMISTIC success contrasts
excluded 0 on the same side at every level; the exception, PREFETCH at L
= 300 ms, includes 0 only at h = 0, where the benefit is removed by
design. All six failure strata included 0 at every level
(\texttt{x4\_interval\_robustness.csv}). This reading is post hoc; it is
a candidate for a future protocol, not a revised verdict.

\textbf{X6, narrow S1 reading.} The frozen robustness rule recomputed
with S1 restricted to Element Timing gives the counts reported in
Section 5.4 (\texttt{x6\_narrow\_s1\_reading.csv}).

\textbf{X5, arm order.} Splitting seeds by whether the arm ran before
BLANK in its block (5 to 15 seeds per group) gave differences of group
medians of at most 9.1 ms for contrasts of 85 to 999 ms, at most 3.8 ms
for the failure strata, and 0.05 ms for SKELETON at L = 50 ms
(\texttt{x5\_arm\_order.csv}). No order effect large enough to change a
verdict is visible, but the groups are small and this is not a test.

\subsection{5.6 Expected versus
observed}\label{expected-versus-observed}

Table 4. Pre-stated expectations, observations and the verdict under the
frozen rule.

\begin{longtable}[]{@{}
  >{\raggedright\arraybackslash}p{(\linewidth - 6\tabcolsep) * \real{0.2500}}
  >{\raggedright\arraybackslash}p{(\linewidth - 6\tabcolsep) * \real{0.2500}}
  >{\raggedright\arraybackslash}p{(\linewidth - 6\tabcolsep) * \real{0.2500}}
  >{\raggedright\arraybackslash}p{(\linewidth - 6\tabcolsep) * \real{0.2500}}@{}}
\toprule\noalign{}
\begin{minipage}[b]{\linewidth}\raggedright
Experiment
\end{minipage} & \begin{minipage}[b]{\linewidth}\raggedright
Prediction (frozen)
\end{minipage} & \begin{minipage}[b]{\linewidth}\raggedright
Observation
\end{minipage} & \begin{minipage}[b]{\linewidth}\raggedright
Verdict under the frozen rule
\end{minipage} \\
\midrule\noalign{}
\endhead
\bottomrule\noalign{}
\endlastfoot
EXP-001 & C1 to C4 pass (instrumentation checks) & Pass in both
exploratory runs & Instrumentation validated (exploratory) \\
EXP-002 & PREFETCH about -min(h, L), within 17 ms & Deviations 0.0 to
15.1 ms & No contradiction; H0 rejected at all 4 delays \\
EXP-002 & SKELETON equivalent within 17 ms & Equivalent at L of 100 to
1000 ms; at 50 ms -16.7 ms, TOST lower bound -17.05 ms & Equivalence
shown at 3 of 4 delays; not shown at 50 ms \\
EXP-002 & SKELETON lowers TTFF & TTFF about 28 ms against about the full
delay & Descriptive; TTFF was not tested \\
EXP-003 & OPTIMISTIC large negative at p = 0 & -85.2, -288.5, -998.6 ms
& H0 rejected in all 9 success strata \\
EXP-003 & Mixed advantage shrinks as p rises & Strictly decreasing at
every delay & No contradiction \\
EXP-003 & Wrong-state exposure grows with failures & False state about
the delay on every failed save & Descriptive, as expected \\
EXP-003 & (no direction pre-stated for failure strata) & -1.6 to 2.7 ms
& H0 not rejected in all 6; no equivalence claim \\
EXP-004 & No reversal of any contrast; SKELETON equivalence stable & 19
of 23 robust; 50 ms SKELETON method-dependent; 3 null failure strata
flagged by the sign rule & Literal contradiction criterion met for 4
contrasts; main interpretation qualified at 50 ms (see Section 6) \\
\end{longtable}

\section{6. Discussion}\label{discussion}

\textbf{What the evidence supports.} Two strategies changed time to
usable content in the direction and of the size that their mechanisms
imply. Prefetching removes the hover lead from the wait, and the
measured benefit matched -min(h, L) within 17 ms at every delay, scaled
with h in S3, and kept its sign under CPU throttling and headed
rendering. An optimistic update removes the whole wait on success, at a
rendering cost of about 27 to 30 ms, and the benefit was robust under
every sensitivity factor. These results check the implementation and
instrumentation as much as the strategies: they confirm that the
benchmark recovers effects that are known by construction.

\textbf{Feedback versus usable content.} The skeleton changed first
feedback without changing time to usable content beyond a fraction of a
frame at delays of 100 ms and more, and this equivalence held under
every method and condition tested. That is the pattern the working title
anticipated. At 50 ms the picture is less clean. Under all three paint
methods the skeleton arm showed content earlier, by between about half a
frame and one frame, and the exploratory decomposition places the
difference in when the frame carrying the content started, not in
rendering cost. Because the equivalence verdict changes with the paint
method, and because a one-frame difference cannot pass a 17 ms margin
when the measured headless frame interval is 18.0 to 18.9 ms, we treat
this result as not robust and do not report it as a benefit. We also do
not claim that it is absent: the evidence supports ``a frame-scale,
method-dependent difference of uncertain mechanism'', nothing stronger.

\textbf{Failure and the cost of optimism.} On failed saves no TTUC
difference between the optimistic update and the pessimistic baseline
was detected (estimates -1.6 to 2.7 ms, intervals reaching 5.4 ms; no
equivalence test was registered), but the optimistic arm displayed a
state the server never held for almost the whole delay. In this
benchmark, the measured cost of an optimistic update is therefore
wrong-state exposure; a slower recovery of a few milliseconds cannot be
excluded. The frozen sign rule flagged three failure strata as method-
or condition-dependent. Their estimates are within 1 ms of zero, not
significant, and their intervals include 0 at every level, so the flags
carry no directional information; they show a limitation of a sign rule
without a tolerance band.

\textbf{The frozen contradiction criterion.} Protocol section 4 states
what would contradict the main interpretation: a SKELETON difference
outside the equivalence margin in either direction, or a reversal of any
contrast under the sensitivity conditions. Read literally, this
criterion is met: equivalence was not shown at 50 ms, and three failure
strata changed sign under some level. We report the main interpretation
as qualified rather than refuted because the 50 ms estimate lies inside
the margin with a method-dependent size, and because the three sign
changes concern estimates within 1 ms of zero whose intervals include 0
at every level. That judgement was made after seeing the data. A reader
who applies the criterion literally should conclude that the frozen test
of the main interpretation was failed at the shortest delay and on the
failure strata.

\textbf{Competing explanations.} The protocols named placeholder
rendering cost, prefetch cache overhead and frame alignment on the read
path, and rollback re-rendering, error-message layout shift and callback
scheduling on the write path. Placeholder rendering cost would make the
skeleton slower; no such cost appeared at any delay, including under CPU
4x. Prefetch overhead is visible only at delays where the residual wait
is zero, as a rendering cost of about 28 ms. Frame alignment is the
explanation the data favour for the 50 ms skeleton effect (X2). Rollback
cost would make failure strata positive; the largest failure estimate
was 2.7 ms with an interval reaching 5.4 ms. Error messages caused no
layout shift (0 in every write cell). Hover-lead jitter and arm order do
not change any verdict (X3, X5). Machine load and time drift were
controlled only by a quiet window and are discussed in Section 7.

\textbf{Practical significance.} In this setting, the size of an
optimistic or prefetch benefit is set by the delay or the lead, not by
the strategy's implementation, and a skeleton is not a way to deliver
content sooner. These are statements about one page in one browser; they
suggest which quantities a practitioner should measure, not how users
will experience them.

\section{7. Limitations and threats to
validity}\label{limitations-and-threats-to-validity}

The full list is maintained in \texttt{research/limitations.md}. The
most important are:

\begin{enumerate}
\def\labelenumi{\arabic{enumi}.}
\tightlist
\item
  \textbf{No perception claims.} All results are objective timings.
  Nothing here measures how fast an interface feels; a perceived-speed
  claim would need a human study under an owner-approved protocol.
\item
  \textbf{One page, one browser version, one machine.} Chrome
  158.0.8127.38 on an Apple M1 Pro, mostly headless. Other engines,
  devices, frameworks and DOM sizes are not covered.
\item
  \textbf{Synthetic latency and scripted input.} Delays come from a
  seeded transport with uniform 10\% jitter; there are no retries,
  timeouts or network effects. The hover is scripted with a fixed lead,
  so the prefetch benefit is bounded by design and says nothing about
  how often real users hover.
\item
  \textbf{Paint-time construct.} The main method measures the second
  animation frame after the commit; the alternatives move some estimates
  by up to about 10 ms (PREFETCH at L = 100 ms: -84.9 ms main, -95.2 ms
  Element Timing). Frame-scale effects are method-dependent (Section
  5.4), the headless frame interval of 18.0 to 18.9 ms exceeds the 17 ms
  margin, and the display refresh rate was not recorded in headed runs.
\item
  \textbf{Pending test reading.} DEC-011 awaits the owner's review: it
  fixes which test decides H0 and which alternative paint methods form
  S1. Both readings are shown where they differ (one contrast in Table
  1; the robustness counts in Section 5.4, 19 of 23 robust against 11
  robust and 9 not assessed).
\item
  \textbf{Sensitivity coverage.} S2 to S4 use 10 seeds and a reduced
  matrix (L in \{300, 1000\} ms, p in \{0, 0.3\}); TTFF sensitivity was
  not computable.
\item
  \textbf{Run conditions.} Runs executed under \texttt{nice\ -n\ 10} in
  a coordinator-arranged quiet window, with load averages of 1.4 to 2.4
  at the starts; background services and an editor were still running.
  The frozen generator's fixed sentence that other agents were running
  is inaccurate for these runs (DEC-012).
\item
  \textbf{Failure strata.} They have 12 or 18 seeds and 26 or 57 trial
  pairs per cell, and no equivalence test was registered for them.
\item
  \textbf{Exploratory analyses.} Section 5.5 is post hoc. It can
  suggest, not confirm.
\end{enumerate}

\section{8. Ethics and responsible
research}\label{ethics-and-responsible-research}

The study uses only a synthetic mock page and a seeded schedule; no
personal data, no external dataset and no human participants were
involved, so no ethics approval was sought or needed, and none is
claimed. Results concern no real or proprietary product. Because the
findings are about timing, not experience, they should not be used to
argue that a strategy improves or harms user perception.

\section{9. Conclusion and future
work}\label{conclusion-and-future-work}

In one mock interface in one browser, intent prefetching and optimistic
updates reduced time to usable content by amounts set by the hover lead
and the delay, and both results were robust to the sensitivity factors
tested. On failed saves no difference in time to the final state was
detected, but the optimistic arm displayed a false state for about the
whole delay. A skeleton placeholder changed first feedback but not time
to usable content at delays of 100 ms and more; at 50 ms a frame-scale
difference appeared whose size depends on the paint-time method, and it
is not robust.

Feasible extensions, each needing a new owner-approved protocol: record
alternative paint times for the first-feedback frame so that TTFF
sensitivity can be computed; use a robustness rule with a tolerance band
or intervals; set the equivalence margin from the measured frame
interval; record the display refresh rate in headed runs; treat renderer
frame-loop state as an explicit factor below 100 ms; add retries and
timeouts to the failure model; and register an equivalence test for the
failure strata.

\section{Data and code availability}\label{data-and-code-availability}

All code, configurations, frozen protocols, raw results with SHA-256
manifests, generated tables and figures are in the project repository:
\texttt{{[}OWNER\ PLACEHOLDER:\ repository\ URL,\ not\ yet\ published{]}}.
Clean-run commands are in each experiment report
(\texttt{research/experiments/EXP-00N/report.md}, section 12). Raw
results: \texttt{results/raw/EXP-002/}, \texttt{results/raw/EXP-003/},
\texttt{results/raw/EXP-004/}; exploratory results:
\texttt{results/exploratory/}. The code is released under the SCROLL
Source-Available Non-Commercial Research License v1.0
(\texttt{LICENSE}).

\section{Funding, competing interests and AI
assistance}\label{funding-competing-interests-and-ai-assistance}

\textbf{Funding.} \texttt{{[}OWNER\ PLACEHOLDER:\ funding\ statement{]}}

\textbf{Competing interests.}
\texttt{{[}OWNER\ PLACEHOLDER:\ competing-interest\ statement{]}}

\textbf{AI assistance.} An AI system (Claude, Anthropic) was used as a
research and development assistant: it drafted the protocols for the
owner's approval, wrote the code and tests, ran the experiments under
the frozen protocol, computed the analyses and drafted this manuscript.
It is not an author. Human verification: the owner approved the protocol
freeze, the amendments and the decisions recorded in
\texttt{research/decision\_log.md}; the owner's review of this
manuscript, of DEC-011 and of the placeholders above is pending. An
internal AI review pass
(\texttt{research/publication/review\_pass\_2.md}) is quality control,
not peer review.

\section{References}\label{references}

\phantomsection\label{refs}
\begin{CSLReferences}{1}{0}
\bibitem[\citeproctext]{ref-mdn_lcp}
MDN Web Docs. n.d. {``Largest Contentful Paint {API}.''} Mozilla
Developer Network.
\url{https://developer.mozilla.org/en-US/docs/Web/API/Largest_Contentful_Paint_API}.

\bibitem[\citeproctext]{ref-mejtoft2018}
Mejtoft, Thomas, Arvid Långström, and Ulrik Söderström. 2018. {``The
Effect of Skeleton Screens: Users' Perception of Speed and Ease of
Navigation.''} In \emph{Proceedings of the 36th European Conference on
Cognitive Ergonomics}, 1--4. ACM.
\url{https://doi.org/10.1145/3232078.3232086}.

\bibitem[\citeproctext]{ref-mihajlija_cls}
Mihajlija, Milica, and Philip Walton. n.d. {``Cumulative Layout Shift
({CLS}).''} web.dev. \url{https://web.dev/articles/cls}.

\bibitem[\citeproctext]{ref-w3c_event_timing}
Mocny, Michal. n.d.a. {``Event Timing {API}.''} W3C Working Draft, 13
November 2026. \url{https://www.w3.org/TR/event-timing/}.

\bibitem[\citeproctext]{ref-wicg_layout_instability}
---------. n.d.b. {``Layout Instability {API}.''} WICG Draft Community
Group Report, 8 November 2026.
\url{https://wicg.github.io/layout-instability/}.

\bibitem[\citeproctext]{ref-nielsen1993response}
Nielsen, Jakob. 1993. {``Response Times: The 3 Important Limits.''}
Nielsen Norman Group (excerpt of Usability Engineering, ch. 5).
\url{https://www.nngroup.com/articles/response-times-3-important-limits/}.

\bibitem[\citeproctext]{ref-osmani_optimize_cls}
Osmani, Addy, and Barry Pollard. n.d. {``Optimize Cumulative Layout
Shift.''} web.dev, published 2020-05-05, updated 2025-02-07.
\url{https://web.dev/articles/optimize-cls}.

\bibitem[\citeproctext]{ref-react_useoptimistic}
React. n.d. {``useOptimistic.''} React documentation.
\url{https://react.dev/reference/react/useOptimistic}.

\bibitem[\citeproctext]{ref-tankala_skeleton}
Tankala, Samhita. n.d. {``Skeleton Screens 101.''} Nielsen Norman Group,
4 June 2023, reviewed 2 September 2026.
\url{https://www.nngroup.com/articles/skeleton-screens/}.

\bibitem[\citeproctext]{ref-wagner_inp}
Wagner, Jeremy, and Barry Pollard. n.d. {``Interaction to Next Paint
({INP}).''} web.dev, published 2022-05-06, updated 2025-09-02.
\url{https://web.dev/articles/inp}.

\end{CSLReferences}

The reference list is generated from
\texttt{research/paper/references.bib}. Every entry was opened as a
model-written page summary on 2026-11-24, not as full text; mejtoft2018
is metadata only and is not used for any finding.

\section{Supplementary appendix}\label{supplementary-appendix}

See \texttt{research/paper/supplementary.md}: complete experimental
matrix, configurations and hashes, extended tables, robustness results,
report links and deviations.

\end{document}
